\section{Abtasttheorem}
\label{sec:abttastheorem}

Die folgenden Veranschaulichungen und mathematischen Definitionen orientieren sich an den Darstellungen von
Lyons \cite[Kap. 1]{lyons2011} und Smith \cite[Kap. 3]{smith1997}.

Aus physikalischer Sicht sind Töne Luftdruckschwankungen, die sich in Form von Wellen durch ein Medium, wie
bspw. der Erdatmosphäre ausbreiten. Solche Schwankungen lassen sich mithilfe von Wandlern wie Mikrofonen in
stetige Wechselstromsignale übertragen:

\cite[Kap. 1]{lyons2011}:
\begin{equation}
    y(t) = \sin(2\pi \cdot f \cdot t)
    \label{eq:stetige_gleichung}
\end{equation}

Problematisch hierbei ist, dass ein solches analoges Signal aufgrund seiner Stetigkeit eine überabzählbare
Anzahl an potenziellen Messpunkten aufweist und sich dadurch unmöglich in Form einer Datenstruktur speichern lässt.

Um dieses Problem nun zu lösen und das Signal dennoch zu digitalisieren,
wird eine Abtastung (engl.: \textit{sampling})
in zeitlich gleichmäßigen Abständen vorgenommen.